\section{Miscellany}
\label{section-miscellany}

\noindent
The proofs in this section should not refer to any results except
those from the section on basic notions, Section \ref{section-rings-basic}.

\begin{lemma}
\label{lemma-product-ideals-in-prime}
Let $R$ be a ring, $I$ and $J$ two ideals and $\mathfrak p$ a prime ideal
containing the product $IJ$. Then $\mathfrak{p}$ contains $I$ or $J$.
\end{lemma}

\begin{proof}
Assume the contrary and take $x \in I \setminus \mathfrak p$ and
$y \in J \setminus \mathfrak p$. Their product is an element of
$IJ \subset \mathfrak p$, which contradicts the assumption that
$\mathfrak p$ was prime.
\end{proof}

\begin{lemma}[Prime avoidance]
\label{lemma-silly}
\begin{slogan}
1. In an affine scheme if a finite number of points are contained in an
open subset then they are contained in a smaller principal open subset.
2. Affine opens are cofinal among the neighborhoods of a given finite set
of an affine scheme
\end{slogan}
Let $R$ be a ring. Let $I_i \subset R$, $i = 1, \ldots, r$,
and $J \subset R$ be ideals. Assume
\begin{enumerate}
\item $J \not\subset I_i$ for $i = 1, \ldots, r$, and
\item all but two of $I_i$ are prime ideals.
\end{enumerate}
Then there exists an $x \in J$, $x\not\in I_i$ for all $i$.
\end{lemma}

\begin{proof}
The result is true for $r = 1$. If $r = 2$, then let $x, y \in J$ with
$x \not \in I_1$ and $y \not \in I_2$. We are done unless $x \in I_2$
and $y \in I_1$. Then the element $x + y$ cannot be in $I_1$ (since that
would mean $x + y - y \in I_1$) and it also cannot be in $I_2$.

\medskip\noindent
For $r \geq 3$, assume the result holds for $r - 1$. After renumbering
we may assume that $I_r$ is prime. We may also assume there are no
inclusions among the $I_i$. Pick $x \in J$, $x \not \in I_i$ for all
$i = 1, \ldots, r - 1$. If $x \not\in I_r$ we are done. So assume
$x \in I_r$. If $J I_1 \ldots I_{r - 1} \subset I_r$ then
$J \subset I_r$ (by Lemma \ref{lemma-product-ideals-in-prime}) a contradiction.
Pick $y \in J I_1 \ldots I_{r - 1}$, $y \not \in I_r$. Then $x + y$ works.
\end{proof}

\begin{lemma}
\label{lemma-silly-silly}
Let $R$ be a ring. Let $x \in R$, $I \subset R$ an ideal, and
$\mathfrak p_i$, $i = 1, \ldots, r$ be prime ideals.
Suppose that $x + I \not \subset \mathfrak p_i$ for
$i = 1, \ldots, r$. Then there exists a $y \in I$
such that $x + y \not \in \mathfrak p_i$ for all $i$.
\end{lemma}

\begin{proof}
We may assume there are no inclusions among the $\mathfrak p_i$.
After reordering we may assume $x \not \in \mathfrak p_i$ for $i < s$
and $x \in \mathfrak p_i$ for $i \geq s$. If $s = r + 1$ then we are done.
If not, then we can find $y \in I$ with $y \not \in \mathfrak p_s$.
Choose $f \in \bigcap_{i < s} \mathfrak p_i$ with $f \not \in \mathfrak p_s$.
Then $x + fy$ is not contained in $\mathfrak p_1, \ldots, \mathfrak p_s$.
Thus we win by induction on $s$.
\end{proof}

\begin{lemma}[Chinese remainder]
\label{lemma-chinese-remainder}
Let $R$ be a ring.
\begin{enumerate}
\item If $I_1, \ldots, I_r$ are ideals such that $I_a + I_b = R$
when $a \not = b$, then $I_1 \cap \ldots \cap I_r =
I_1I_2\ldots I_r$ and $R/(I_1I_2\ldots I_r)
\cong R/I_1 \times \ldots \times R/I_r$.
\item If $\mathfrak m_1, \ldots, \mathfrak m_r$ are pairwise distinct maximal
ideals then $\mathfrak m_a + \mathfrak m_b = R$ for $a \not = b$ and the
above applies.
\end{enumerate}
\end{lemma}

\begin{proof}
Let us first prove $I_1 \cap \ldots \cap I_r = I_1 \ldots I_r$
as this will also imply the injectivity of the induced ring
homomorphism $R/(I_1 \ldots I_r) \rightarrow R/I_1 \times \ldots \times R/I_r$.
The inclusion $I_1 \cap \ldots \cap I_r \supset I_1 \ldots I_r$ is always
fulfilled since ideals are closed under multiplication with arbitrary ring
elements. To prove the other inclusion, we claim that the ideals
$$
I_1 \ldots \hat I_i \ldots I_r,\quad i = 1, \ldots, r
$$
generate the ring $R$. We prove this by induction on $r$. It holds when
$r = 2$. If $r > 2$, then we see that $R$ is the sum of the ideals
$I_1 \ldots \hat I_i \ldots I_{r - 1}$, $i = 1, \ldots, r - 1$.
Hence $I_r$ is the sum of the ideals
$I_1 \ldots \hat I_i \ldots I_r$, $i = 1, \ldots, r - 1$.
Applying the same argument with the reverse ordering on the ideals
we see that $I_1$ is the sum of the ideals
$I_1 \ldots \hat I_i \ldots I_r$, $i = 2, \ldots, r$.
Since $R = I_1 + I_r$ by assumption we see that $R$ is the sum of the
ideals displayed above. Therefore we can find elements
$a_i \in I_1 \ldots \hat I_i \ldots I_r$
such that their sum is one. Multiplying this equation by an element
of $I_1 \cap \ldots \cap I_r$ gives the other inclusion.
It remains to show that the canonical map
$R/(I_1 \ldots I_r) \rightarrow R/I_1 \times \ldots \times R/I_r$
is surjective. For this, consider its action on the equation
$1 = \sum_{i=1}^r a_i$ we derived above. On the one hand, a
ring morphism sends 1 to 1 and on the other hand, the image of any
$a_i$ is zero in $R/I_j$ for $j \neq i$. Therefore, the image of $a_i$
in $R/I_i$ is the identity. So given any element
$(\bar{b_1}, \ldots, \bar{b_r}) \in R/I_1 \times \ldots \times R/I_r$,
the element $\sum_{i=1}^r a_i \cdot b_i$ is an inverse image in $R$.

\medskip\noindent
To see (2), by the very definition of being distinct maximal ideals, we have
$\mathfrak{m}_a + \mathfrak{m}_b = R$ for $a \neq b$ and so the above applies.
\end{proof}

\begin{lemma}
\label{lemma-matrix-left-inverse}
Let $R$ be a ring. Let $n \geq m$. Let $A$ be an
$n \times m$ matrix with coefficients in $R$. Let $J \subset R$
be the ideal generated by the $m \times m$ minors of $A$.
\begin{enumerate}
\item For any $f \in J$ there exists a $m \times n$ matrix $B$
such that $BA = f 1_{m \times m}$.
\item If $f \in R$ and $BA = f 1_{m \times m}$ for some $m \times n$ matrix
$B$, then $f^m \in J$.
\end{enumerate}
\end{lemma}

\begin{proof}
For $I \subset \{1, \ldots, n\}$ with $|I| = m$, we denote
by $E_I$ the $m \times n$ matrix of the projection
$$
R^{\oplus n} = \bigoplus\nolimits_{i \in \{1, \ldots, n\}} R
\longrightarrow \bigoplus\nolimits_{i \in I} R
$$
and set $A_I = E_I A$, i.e., $A_I$ is the $m \times m$ matrix
whose rows are the rows of $A$ with indices in $I$.
Let $B_I$ be the adjugate (transpose of
cofactor) matrix to $A_I$, i.e., such that
$A_I B_I = B_I A_I = \det(A_I) 1_{m \times m}$.
The $m \times m$ minors of $A$ are the determinants $\det A_I$
for all the $I \subset \{1, \ldots, n\}$ with $|I| = m$.
If $f \in J$ then we can write $f = \sum c_I \det(A_I)$ for some
$c_I \in R$. Set $B = \sum c_I B_I E_I$ to see that (1) holds.

\medskip\noindent
If $f 1_{m \times m} = BA$ then by the
Cauchy-Binet formula (\ref{item-cauchy-binet}) we
have $f^m = \sum b_I \det(A_I)$ where $b_I$ is the determinant
of the $m \times m$ matrix whose columns are the columns of $B$ with
indices in $I$.
\end{proof}

\begin{lemma}
\label{lemma-matrix-right-inverse}
Let $R$ be a ring. Let $n \geq m$. Let $A = (a_{ij})$ be an
$n \times m$ matrix with coefficients in $R$, written in block form
as
$$
A =
\left(
\begin{matrix}
A_1 \\
A_2
\end{matrix}
\right)
$$
where $A_1$ has size $m \times m$. Let $B$ be the adjugate (transpose of
cofactor) matrix to $A_1$. Then
$$
AB = 
\left(
\begin{matrix}
f 1_{m \times m} \\
C
\end{matrix}
\right)
$$
where $f = \det(A_1)$ and $c_{ij}$ is (up to sign) the determinant of the
$m \times m$ minor of $A$ corresponding to the rows
$1, \ldots, \hat j, \ldots, m, i$.
\end{lemma}

\begin{proof}
Since the adjugate has the property $A_1B = B A_1 = f$ the first block
of the expression for $AB$ is correct. Note that
$$
c_{ij} = \sum\nolimits_k a_{ik}b_{kj} = \sum (-1)^{j + k}a_{ik} \det(A_1^{jk})
$$
where $A_1^{ij}$ means $A_1$ with the $j$th row and $k$th column removed.
This last expression is the row expansion of the determinant of the matrix
in the statement of the lemma.
\end{proof}

\begin{lemma}
\label{lemma-map-cannot-be-injective}
\begin{slogan}
A map of finite free modules cannot be injective if the source has
rank bigger than the target.
\end{slogan}
Let $R$ be a nonzero ring. Let $n \geq 1$. Let $M$ be an $R$-module generated
by $< n$ elements. Then any $R$-module map $f : R^{\oplus n} \to M$ has a
nonzero kernel.
\end{lemma}

\begin{proof}
Choose a surjection $R^{\oplus n - 1} \to M$.
We may lift the map $f$ to a map $f' : R^{\oplus n} \to R^{\oplus n - 1}$
(Lemma \ref{lemma-lift-map}).
It suffices to prove $f'$ has a nonzero kernel.
The map $f' : R^{\oplus n} \to R^{\oplus n - 1}$ is given by a
matrix $A = (a_{ij})$. If one of the $a_{ij}$ is not nilpotent, say
$a = a_{ij}$ is not, then we can replace $R$ by the localization $R_a$
and we may assume $a_{ij}$ is a unit. Since if we find a nonzero kernel
after localization then there was a nonzero kernel to start with as
localization is exact, see Proposition \ref{proposition-localization-exact}.
In this case we can do a base change on both $R^{\oplus n}$
and $R^{\oplus n - 1}$ and reduce to the case where
$$
A =
\left(
\begin{matrix}
1 & 0 & 0 & \ldots \\
0 & a_{22} & a_{23} & \ldots \\
0 & a_{32} & \ldots \\
\ldots & \ldots
\end{matrix}
\right)
$$
Hence in this case we win by induction on $n$. If not then each
$a_{ij}$ is nilpotent. Set $I = (a_{ij}) \subset R$. Note that
$I^{m + 1} = 0$ for some $m \geq 0$. Let $m$ be the largest integer
such that $I^m \not = 0$. Then we see that $(I^m)^{\oplus n}$ is
contained in the kernel of the map and we win.
\end{proof}

\begin{lemma}
\label{lemma-rank}
\begin{slogan}
The rank of a finite free module is well defined.
\end{slogan}
Let $R$ be a nonzero ring. Let $n, m \geq 0$ be integers.
If $R^{\oplus n}$ is isomorphic to $R^{\oplus m}$ as
$R$-modules, then $n = m$.
\end{lemma}

\begin{proof}
Immediate from Lemma \ref{lemma-map-cannot-be-injective}.
\end{proof}





